% eLight Journal Template (Springer Nature / Light Publishing Group)
% Publisher: SpringerOpen (open access)
% Style: Springer Nature sn-jnl class
% Abstract: ≤200 words; no citations, no abbreviations
% Keywords: 3–10 keywords
% Figures: ≥300 dpi; multi-panel combined into one file; uploaded separately
% References: numbered, [1] style
% Double-blind review: remove author info for initial submission

\documentclass[lineno]{sn-jnl}
% Options: lineno (line numbers), iicol (two-column), referee (double-space)
% Download sn-jnl.cls from: https://www.springernature.com/gp/authors/campaigns/latex-author-support

\usepackage{amsmath,amssymb,bm}
\usepackage{graphicx}
\usepackage{xcolor}
\usepackage{hyperref}
\usepackage{siunitx}
\usepackage{booktabs}
\usepackage[version=4]{mhchem}

% Placeholder macro for missing figures
\newcommand{\placeholder}[1]{%
  \begin{center}
    \fbox{\parbox{0.9\columnwidth}{\centering\small\textit{[FIGURE: #1]}}}
  \end{center}
}

% ─── Title ───────────────────────────────────────────────────────────────────
\title[Short title (≤40 chars)]{Full Title of the eLight Manuscript (7–25 words)}

% ─── Authors (remove for double-blind initial submission) ─────────────────────
\author[1]{\fnm{First} \sur{Author}}
\author[1,2]{\fnm{Second} \sur{Author}}
\author*[1,2]{\fnm{Corresponding} \sur{Author}}

\affil[1]{\orgdiv{Department of Physics}, \orgname{University Name}, \orgaddress{\city{City}, \country{Country}}}
\affil[2]{\orgdiv{State Key Laboratory of ...}, \orgname{Institute Name}, \orgaddress{\city{City}, \country{Country}}}

\corresp{corresponding@email.com}

\begin{document}

\maketitle

% ─── Abstract ────────────────────────────────────────────────────────────────
% ≤200 words; no citations; minimize abbreviations
% Structure: background → problem → approach → result → significance
\begin{abstract}
Background sentence placing the work in context of modern photonics/light science.
Problem sentence identifying the unresolved challenge or limitation of existing approaches.
``Here, we [demonstrate/show/report]...'' — approach sentence.
Key quantitative result (e.g., ``achieving a Q-factor of X, representing a Y-fold improvement'').
Significance sentence on impact for future photonic devices/systems.
\end{abstract}

% ─── Keywords ────────────────────────────────────────────────────────────────
\keywords{photonics, nanophotonics, light–matter interaction, keyword4, keyword5}

% ─── Introduction ────────────────────────────────────────────────────────────
\section{Introduction}\label{sec:intro}

Broad context: light–matter interaction, photonic integration, or relevant subfield.\cite{ref1}
Specific challenge: what physical phenomenon or device functionality remains elusive.\cite{ref2,ref3}
Prior work and why it is insufficient.\cite{ref4,ref5}
``In this work, we...''\cite{ref6}
Paper outline (optional).

% ─── Results ─────────────────────────────────────────────────────────────────
\section{Results}\label{sec:results}

\subsection{Sample Design and Fabrication}

\begin{figure}[htbp]
  \centering
  \placeholder{Device schematic, fabrication process, and structural characterization}
  \caption{\textbf{Device architecture and fabrication.}
           \textbf{a} Schematic of the photonic structure.
           \textbf{b} SEM cross-section.
           \textbf{c} Measured layer thickness profile.
           Scale bars: \SI{1}{\micro\meter}.}
  \label{fig:device}
\end{figure}

\subsection{Optical Characterization}

\begin{figure*}[htbp]
  \centering
  \placeholder{Linear and nonlinear optical spectra, mode profiles, dispersion}
  \caption{\textbf{Optical spectroscopy.}
           \textbf{a}--\textbf{b} Reflection/transmission spectra (experiment vs.\ theory).
           \textbf{c} Photonic band diagram.
           \textbf{d} Near-field map at resonance wavelength.}
  \label{fig:optics}
\end{figure*}

\subsection{Key Result}

\begin{figure}[htbp]
  \centering
  \placeholder{Core device performance: emission, lasing, modulation, or sensing metrics}
  \caption{\textbf{[Key result title].}
           Panel descriptions (a)--(d).}
  \label{fig:result}
\end{figure}

% ─── Discussion ──────────────────────────────────────────────────────────────
\section{Discussion}\label{sec:discussion}

Physical interpretation of the key result.
Comparison with state of the art (table if helpful).
Limitations and future improvements.

% ─── Conclusion ──────────────────────────────────────────────────────────────
\section{Conclusion}\label{sec:conclusion}

In summary, we have demonstrated...
Key finding with numbers.
Outlook for applications and follow-on work.

% ─── Methods ─────────────────────────────────────────────────────────────────
\section{Methods}\label{sec:methods}

\subsection{Device Fabrication}

\subsection{Optical Measurement Setup}

\subsection{Numerical Simulations}

\subsection{Data Analysis and Statistics}

% ─── Declarations ────────────────────────────────────────────────────────────
\section*{Declarations}

\subsection*{Competing interests}
The authors declare that they have no competing interests.

\subsection*{Authors' contributions}
% CRediT: A.O.: conceptualization, investigation, writing—original draft.
% C.A.: supervision, funding acquisition, writing—review and editing.

\subsection*{Funding}
This work was supported by [Agency] (Grant No. XXX).

\subsection*{Availability of data and materials}
The data that support the findings are available from the corresponding author on reasonable request.

% ─── Supplementary Information ───────────────────────────────────────────────
% Upload as separate file; reference in text as (Supplementary Fig. S1)
% \begin{suppinfo}
%   Supplementary Fig. S1: ...
% \end{suppinfo}

% ─── References ──────────────────────────────────────────────────────────────
% Numbered style [1]; use natbib \cite{}
\bibliography{refs}

\end{document}
